% 第 12 章　平衡、材料与流体
\chapter{平衡、材料与流体}

\step{反常识开场}

\begin{mainline}
先做一个你今晚就能做的实验。取一只普通玻璃杯，装满水，水面略高于杯口，盖上一张硬一点的纸牌（扑克牌、名片都行），用手掌压住纸牌，把杯子整个倒过来，然后松开手掌。纸牌没有掉。满满一杯水——大约半斤重——被一张几克重的纸牌托在半空。

水在纸牌上面，纸牌下面只有空气。是水把纸牌“粘”住了吗？是杯口“吸”住了牌吗？如果你真的以为是“吸”，请继续读本章，因为这个“吸”字骗过了几乎所有没学过物理的人。真正托住纸牌的那位大力士，此刻也正压在你的肩膀上，只是你从来感觉不到它。

再看两个同一家族的现象。台风过境时，被掀掉的屋顶往往不是被风\emph{从下往上顶}飞的，而是风贴着屋顶\emph{刮过去}，屋顶却像被一只无形的手向上拔掉。几十吨重的客机靠两片薄薄的机翼在万米高空飞行，机翼上下表面受到的力之差，要稳稳托住几百名乘客和全部燃油。一张纸牌、一片屋顶、一只机翼，它们背后是同一套规律：会流动的东西——水和空气——怎样推、怎样托、怎样掀。

这一章还要回答另一组看似无关的问题：跷跷板上的大人为什么自觉往前坐？高压水枪为什么能把水喷得又急又远？橡皮筋能拉长好几倍再缩回去，粉笔却一掰就断，固体凭什么“硬”、又凭什么“断”？到本章结尾你会发现：平衡、材料、流体这三件事，其实被同一条线索穿着——\emph{力和压强如何在物体内部和物体之间传递与抵消}。
\end{mainline}

\step{先猜一猜}

老规矩，先押注，再看答案。

第一题。倒扣水杯实验中，纸牌不掉的最主要原因是什么？（a）水把纸牌润湿后“粘”住了牌；（b）杯口边缘对纸牌有“吸力”；（c）牌下面的空气在向上顶它。如果你选（c），请再猜：空气向上顶的力大约有多大——几克重、几百克重，还是几十千克重？

第二题。大人七十千克，小孩三十千克，两人玩跷跷板。大人坐在离支点一米处，小孩要坐在离支点多远处才能压平他？（a）也是一米，两边坐对称才公平；（b）比一米远；（c）比一米近。再想一想：如果大人干脆坐\emph{在支点上}，小孩还能把他跷起来吗？

第三题。一颗小铁钉扔进水里立刻沉底，一艘几万吨的钢铁巨轮却稳稳浮着。如果“重的沉、轻的浮”这条经验是对的，问题出在哪？铁钉和巨轮，差的是什么？

第四题。把两根吸管同时含在嘴里，一根插进饮料杯，另一根露在杯子外面的空气中，然后用力吸。你能喝到饮料吗？

写下你的四个答案。本章结束时我们逐一对账。

\step{动手实验}

\begin{experiment}[会潜水的“小药瓶”]
\textbf{材料}：一只带盖的透明软塑料瓶（矿泉水瓶即可）、一只小药瓶或口服液小玻璃瓶（没有的话，用一截折弯的吸管夹上回形针代替）、水。

\textbf{步骤一}：往小药瓶里灌一部分水，让它倒扣着放进一盆水里时\emph{恰好浮着}——瓶内留一个小气泡，瓶口朝下。这个气泡就是全部机关。反复调整水量，直到它浮得很勉强，轻轻一按就沉。我们叫它“沉浮子”。

\textbf{步骤二}：把沉浮子放进灌满水的塑料瓶，拧紧瓶盖。现在它浮在瓶顶。

\textbf{步骤三}：用双手使劲捏瓶身。沉浮子会缓缓下沉；手一松，它又浮上来。捏得越用力，它沉得越深。调好水量的话，你甚至能让它悬在瓶子中间不上不下。

\textbf{记录}：捏瓶身时，你的手并没有碰到沉浮子，是什么“命令”它下沉的？瓶里被压缩的是什么？沉浮子内部的气泡在你捏瓶身时变大了还是变小了？把你的观察画下来——这个实验里藏着本章最重要的三个概念：压强的传递、压缩与浮力。答案在“回头解释开场现象”一节揭晓。
\end{experiment}

还有一个零成本实验，只需一张纸条。撕一条宽约两厘米、长约十五厘米的纸条，用一只手的两根手指捏住一端，让纸条自然下垂，纸条位于你嘴唇\emph{下方}。先猜：沿水平方向用力向纸条\emph{上方}吹气，纸条会被吹得更低，还是升起来？吹一口试试。纸条会向上飘起，几乎升到水平。你的直觉多半会猜错——直觉认为“吹气就是把东西推开”，可纸条偏偏朝着气流的方向凑过去。把这条惹事的纸条留好，本章后半段要为它平反。

\step{旧模型遇到麻烦}

\textbf{旧模型一：“重的沉，轻的浮。”}这句经验听起来天经地义：石头沉、木头浮，铁沉、泡沫浮。可一颗二十克的铁钉沉底，一艘四万吨的铁船却浮着——同样是铁，重的那个反而浮了。有人补救说“船里有空气，平均下来轻了”，这已经接近真理，但它同时承认了旧模型的死刑：决定沉浮的从来不是“这个东西多重”这\emph{一个}数字。旧模型连一个统一的判断标准都给不出。

\textbf{旧模型二：“吸”是一种拉力。}喝饮料时，我们觉得是嘴把饮料“吸”上来的，仿佛嘴伸出一根看不见的绳子把液面往上拽。可这个模型在三个场合接连翻船。其一，把饮料装在完全密封的硬瓶里，插一根吸管，瓶子不漏气，你吸到脸酸也喝不上来——“拉力”还在，饮料不动。其二，就是先猜一猜的第四题：嘴里同时含两根吸管，一根在饮料里、一根在空气里，照样喝不到。其三，最致命的：如果把瓶子上方的空气抽走，“吸力”再强也一滴都上不来。历史上有位更倒霉的受害者——矿井抽水机。几百年前矿工就发现，无论泵造得多精良，从井里抽水最多只能抽上约十米高，再高一厘米都不行。如果“吸”是泵的拉力，力气大就该抽得更高，为什么全欧洲最好的泵都撞在同一堵十米高的墙上？

\textbf{旧模型三：“东西不动，就是没受力”或者“力抵消了”。}前半句的错误本书第 6 章已经拆过：书放在桌上，重力和支持力都作用着，只是合力为零。但“力抵消了就万事大吉”也还不够。看跷跷板：大人小孩各坐一端，如果两人重力恰好相等、支点在中间，板子平衡；把支点往大人那边挪一挪，还是那两个人、还是那两个力，合力依然为零，板子却翻了。合力为零并不能保证一样东西不\emph{转}。可见“不动”是一个比“合力为零”更挑剔的状态，我们漏掉了某笔和转动有关的账。

\textbf{旧模型四：“风吹在物体上就是推它。”}吹纸条实验直接打脸：气流从纸条上方掠过，纸条不向\emph{下}躲，反而向\emph{上}迎。同样离谱的还有：把两张纸并排挂起来，相隔几厘米，向两张纸\emph{中间}吹气——直觉说它们会被吹开，实际上两张纸越贴越近。旧模型把流体的力想成“推一下”这么简单，显然不够用了。

四条麻烦指向同一个缺口：我们从来没有认真定义过“流体怎样对物体施力”。液体和气体不能用手去抓，它们的推和托全靠一样东西——压强。

\step{发明新概念}

\begin{mainline}
\textbf{第一个新概念：力矩——“转动效果”的记账单位。}跷跷板的教训是：同样大的力，作用点离支点越远，掀动转动的本事越大。推门时你绝不会去推门轴旁边，而是推离门轴最远的门把手；拧生锈的螺丝要换长柄扳手。我们把“力”乘以“力臂”（支点到力的作用线的垂直距离）叫做\emph{力矩}。力矩大，改变转动的本事就大。于是一个物体真正“不动”需要两个条件同时成立：所有的力加总为零（不加速平移），所有的力矩加总为零（不加速转动）。前者管“走不走”，后者管“转不转”。这就是为什么大人要往前坐——他减小的不是体重，是自己的力臂。

\textbf{第二个新概念：压强——摊在面积上的力。}同样用五十牛的力，用手掌推墙，墙纹丝不动；用图钉的针尖推墙，针尖轻松扎进去。力没有变，变的是它被\emph{分摊}到的面积。我们把“单位面积上分担的力”叫做\emph{压强}。它像什么？像把一笔账摊到很多份上：一百元的账单十个人分，每人十元；压强就是把一份力摊到面积上，面积越小，每一点分到的越多。它不像什么？它不像力本身——力是总账，压强是人均账，说“这个物体受到多大的压强”和“受到多大的力”是两回事。图钉的全部秘密就在这一笔账的两头：钉帽面积大，手指受的压强小，不疼；针尖面积不到钉帽的万分之一，同样的力集中成上万倍的压强，木头招架不住。坦克用履带、滑雪板做得又宽又长，是同一原理的反方向应用。

\textbf{第三个新概念：流体内部的压强——越深处，头顶的货越多。}水下的任何一点，都在替它头顶上那整条水柱扛着重量。潜得越深，扛的水柱越长，压强越大。这就是为什么水坝要修成下宽上窄的梯形——坝底承受的压强比水面大得多。更妙的是，液体在某处的压强是\emph{向四面八方}一样大的：它向左向右向下向上都一样挤。密闭的流体还有一条重要脾气：你在某一处给它加多大压强，它就把这份压强\emph{原封不动地传到每一个角落}——这叫帕斯卡原理。汽车刹车就是靠它：脚踩刹车踏板，推动一个小活塞给刹车油加压，这份压强顺着油管传到四个轮子上更大的活塞，把刹车片死死压到轮盘上。

\textbf{第四个新概念：大气压与浮力——我们住在空气海洋的海底。}空气看不见，但它有重量。地球表面每一寸土地上方都压着一根从地面直通太空边缘的空气柱。这根柱子摊到每平方米地面上的重量，大约是\emph{十吨}。你不觉得被压垮，是因为你身体内部、细胞里、肺里的压强和外面一模一样大，内外抵消——就像深海鱼不觉得水压，因为里外平衡。而倒扣水杯实验里，杯内水面那点压强远小于杯外空气从下面顶住纸牌的压强，净效果是空气把牌和水一起托住。纸牌不是被“吸”住的，是被从下面\emph{顶}住的。

一旦接受“流体内部压强随深度增大”，浮力就自动蹦出来了。一个浸在水中的方块：上表面浅，受到的向下压强小；下表面深，受到的向上压强大；前后左右互相抵消。上下一减，水对方块的净效果是一个向上的推力，这就是浮力。这个推力多大？阿基米德泡进浴缸时发现：恰好等于这个方块\emph{排开的那部分水}所受的重力。注意这个表述的狡猾之处：它根本不提方块自己多重、是什么材料——只提“如果这个位置装的是水，那些水有多重”。铁钉排开二十克水的体积，浮力只有二十克水的重量，托不住；巨轮造成空心的大肚子，排开几万吨水的体积，浮力就有几万吨。沉浮的裁判从来不是“你多重”，而是“你的体积能换来多大的浮力，够不够抵消你的重量”。

\textbf{第五个新概念：流进多少，流出多少；加速的水，暴露了压强差。}水在管子里流动时，如果不可压缩，那么每一秒钟流进一段管子的水量，必须等于流出的水量——水不会凭空消失，也塞不下更多。管子变细的地方，水只能流得更快。捏扁浇水软管的头，水射得又急又远，就是这个道理。再看第二张王牌：一段水从粗管冲进细管，速度变快了——\emph{它加速了}。而第 6 章告诉我们，加速需要合力。谁来推这段水？只能是它身后的压强比身前大。于是我们反推出一条反直觉的结论：\emph{稳定流动中，流得快的地方，压强反而低}。这就是伯努利关系的物理内核：不是“快”本身神秘地降低了压强，而是压强差把水推着加速，快是\emph{结果}，低压强是\emph{原因}——这笔账的正式算法是能量，我们马上翻译。

\textbf{第六个新概念：弹性——固体都是藏起来的弹簧。}拉一根橡皮筋，拉力越大，伸得越长，松手就弹回去。弹簧秤利用的正是“伸长量和拉力成正比”这条规律（胡克定律）。深刻的地方在于：\emph{任何}固体都是弹簧，只是硬得离谱。你坐的椅子被你压弯了大约一根头发丝的量，再以同样大的力把你弹回来；钢桥在卡车驶过时向下弯几毫米再弹回。固体能这样，是因为原子之间既不肯被挤近、也不肯被拉远，像亿万个微型弹簧连在一起。这个比喻像的部分是：微小的挤压和拉伸确实近似正比于恢复力。它不像的部分是：原子之间并没有真的弹簧，只有电磁相互作用形成的“近了排斥、远了吸引”的平衡位置（第 54 章会细讲），而且拉得足够远时，这种“弹簧”会永久失灵。一根材料能承受多大的力而不坏，叫强度；能吸收多少损伤而不断，叫韧性。粉笔一掰就断，钢丝能吊起汽车——差异不在“硬不硬”，而在内部原子之间那笔弹簧账能记到什么程度。
\end{mainline}

\step{一句话模型}

\begin{oneline}
平衡不是“没有力”，而是所有的推和拉互相抵消、所有的转动效果也互相抵消；流体对物体的力全来自压强差——水下物体被深处更大的压强向上顶起，流动的水被压强差推着加速，而沉浮的裁判只有一句话：你排开的那份水有多重。
\end{oneline}

\step{数学翻译器}

\begin{mainline}
\textbf{力矩与平衡。}设力为 \(F\)，力臂（支点到力的作用线的垂直距离）为 \(L\)，力矩记作
\[
  \tau = F\,L .
\]
它描述什么：这个力让物体绕支点转动的效果有多强。为什么这样写：实验上，转动效果既正比于力、也正比于力臂，二者是相乘关系——力加倍或力臂加倍，效果都加倍。平衡条件是两行：
\[
  F_{\text{合}} = 0, \qquad \tau_{\text{合}} = 0 .
\]
例行检查。令 \(L=0\)：力正对着支点作用，力矩为零——再大力气推门轴，门也不转，与经验一致。方向反转：让同一个力从“向上抬”换成“向下压”，力矩从逆时针变成顺时针，符号变号，说明平衡是顺时针力矩与逆时针力矩的\emph{对消}。算一笔跷跷板的账：大人 \(70\) 千克坐在支点一侧一米处，力矩大小为 \(70g\times 1\) 米；小孩 \(30\) 千克要产生同样大的反向力矩，需坐在另一侧 \(70\times 1/30\approx 2.3\) 米处——果然是“比一米远”，第二题选（b）。而大人若坐在支点上，\(L=0\)，力矩恒为零，小孩怎样用力也跷不动他。

\textbf{压强。}
\[
  p = \frac{F}{S} .
\]
它描述什么：力在面积 \(S\) 上摊得多“集中”。检查：力不变、面积扩大十倍，压强降为十分之一——履带和滑雪板成立；面积趋于极小，压强极大——针尖成立。算真实数据：你用 \(50\) 牛的力按图钉，钉帽面积约 \(1\) 平方厘米，手指承受 \(50/10^{-4}=5\times10^{5}\) 帕，毫无痛感；钉尖面积约 \(0.01\) 平方毫米，压强高达 \(50/10^{-8}=5\times10^{9}\) 帕——一万倍的差距，同一根手指施加的同一笔力。

\textbf{液体压强与大气压。}水面下深度 \(h\) 处，
\[
  p = p_{0} + \rho g h,
\]
其中 \(p_{0}\) 是液面上的大气压（约 \(1.0\times10^{5}\) 帕），\(\rho\) 是水的密度（每立方米 \(1000\) 千克），\(g\) 取约 \(10\) 牛每千克。为什么这样写：该点要扛起头顶大气和一段高 \(h\)、底面积单位的水柱，水柱重量就是 \(\rho g h\)。检查：令 \(h=0\)，退化为 \(p=p_{0}\)，液面处只剩大气压，正确。代入 \(h=10\) 米：\(\rho g h=1000\times10\times10=10^{5}\) 帕，恰好又是一个大气压——所以每下潜十米，相当于多扛一整层大气。欧洲抽水机的“十米之墙”到此破案：抽水机能做的只是让管内气压接近零，真正把水推上来的是井面的大气压，而一个大气压最多只推得动约十米高的水柱——泵的“吸力”再强也无济于事，因为干活的一直不是它。

\textbf{浮力。}阿基米德原理：
\[
  F_{\text{浮}} = \rho_{\text{液}}\, g\, V_{\text{排}} .
\]
它描述什么：物体排开体积 \(V_{\text{排}}\) 的流体后，受到的净向上推力。检查三发。令 \(V_{\text{排}}=0\)：不接触流体，浮力为零，平凡但必要。物体全部没入水中后，\(V_{\text{排}}\) 不再随深度变化——所以完全浸没后浮力与深度无关，石头下沉过程中浮力不增不减。令物体密度恰等于水的密度：它的重力等于同体积水的重力，浮力恰好托住，悬浮——鱼靠鱼鳔微调体积干的就是这件事。算真人：一个 \(60\) 千克的人体积约 \(60\) 升，完全浸没时浮力约等于 \(60\) 千克水的重力，与体重几乎打平——所以深吸一口气（胸廓变大、\(V_{\text{排}}\) 变大）就能浮起来，呼气就沉，你早就用身体验证过这条公式。

\textbf{连续性方程。}不可压缩流体稳定流过一根管子，截面 \(S\) 与流速 \(v\) 满足
\[
  S_{1}v_{1} = S_{2}v_{2} .
\]
它描述什么：每秒通过任何截面的水量（流量）都相等——流进多少，流出多少。检查：截面减半，流速加倍，捏水管射得远，成立；让 \(S_{2}\) 趋于极大（水管通进水库），流速趋于零，大河入海口水面平缓，也成立。

\textbf{伯努利关系。}沿同一条流线，
\[
  p + \tfrac{1}{2}\rho v^{2} + \rho g h = \text{常量}.
\]
读法：\(\tfrac{1}{2}\rho v^{2}\) 是单位体积水的动能，\(\rho g h\) 是单位体积水的重力势能，\(p\) 是单位体积水“携带着的压强能”——三者之和沿流线守恒，这是能量守恒（第 9 章）穿在流水身上的外衣。检查：令流速处处相等，两项相消，剩下 \(p+\rho g h=\) 常量——正是静水压强公式穿了个马甲回来，自洽。令高度不变：流速大的地方压强必小，吹纸条实验的判决书落地——纸条上方气流快、压强低，下方静止空气压强高，纸条被从下面顶起来。
\end{mainline}

\begin{center}
\begin{tikzpicture}[scale=1.0]
  % 支点
  \draw[thick] (0,0) -- (-0.45,-0.7) -- (0.45,-0.7) -- cycle;
  \fill[pattern=north east lines] (-0.9,-1.0) rectangle (0.9,-0.7);
  \node[below] at (0,-1.0) {支点};
  % 跷跷板
  \draw[very thick] (-3.6,0.3) -- (3.0,-0.25);
  % 小孩（左，远）与大人（右，近）
  \draw[thick,fill=blue!10] (-3.35,0.29) circle (0.25);
  \node[above=2pt] at (-3.35,0.55) {小孩};
  \draw[-{Stealth[length=2.5mm]},thick,blue!70!black] (-3.35,0.29) -- (-3.35,-0.71) node[midway,left=1pt] {\(30g\)};
  \draw[thick,fill=red!10] (1.3,-0.14) circle (0.3);
  \node[above=2pt] at (1.3,0.22) {大人};
  \draw[-{Stealth[length=2.5mm]},thick,red!70!black] (1.3,-0.14) -- (1.3,-1.44) node[midway,right=1pt] {\(70g\)};
  % 力臂标注
  \draw[dashed] (0,0) -- (-3.35,0.28);
  \node[above] at (-1.7,0.2) {力臂约 \(2.3\) 米};
  \draw[dashed] (0,0) -- (1.3,-0.15);
  \node[below] at (0.75,0.05) {力臂 \(1\) 米};
\end{tikzpicture}
\end{center}

\begin{center}
\begin{tikzpicture}[scale=1.0]
  % 水
  \fill[blue!6] (0,0) rectangle (7,-3.2);
  \draw[thick,blue!50!black] (0,0) -- (7,0);
  \node[above left] at (0,0) {水面};
  % 方块
  \draw[thick,fill=gray!20] (2.6,-1.2) rectangle (4.4,-2.6);
  \node at (3.5,-1.9) {浸没的方块};
  % 上表面压强（小）
  \draw[-{Stealth[length=2.5mm]},thick,red!70!black] (3.2,-0.35) -- (3.2,-1.15);
  \draw[-{Stealth[length=2.5mm]},thick,red!70!black] (3.8,-0.35) -- (3.8,-1.15);
  \node[above] at (3.5,-0.3) {上表面：向下的压强较小};
  % 下表面压强（大）
  \draw[-{Stealth[length=3.5mm]},very thick,blue!70!black] (3.2,-3.35) -- (3.2,-2.65);
  \draw[-{Stealth[length=3.5mm]},very thick,blue!70!black] (3.8,-3.35) -- (3.8,-2.65);
  \node[below] at (3.5,-3.4) {下表面更深：向上的压强较大};
  % 侧面抵消
  \draw[-{Stealth[length=2.5mm]},thick] (1.6,-1.9) -- (2.55,-1.9);
  \draw[-{Stealth[length=2.5mm]},thick] (5.4,-1.9) -- (4.45,-1.9);
  \node[right] at (5.4,-1.9) {侧面互相抵消};
  % 净浮力
  \draw[-{Stealth[length=4mm]},very thick,orange!80!black] (5.6,-2.6) -- (5.6,-0.9);
  \node[right] at (5.6,-1.75) {净效果：浮力向上};
\end{tikzpicture}
\end{center}

\begin{mathbridge}
伯努利关系可以直接从功能关系推出。取一小段体积为 \(\Delta V\) 的水，从截面 1 流到截面 2。身后的水对它做正功 \(p_{1}\Delta V\)，身前的水对它做负功 \(p_{2}\Delta V\)，合功为 \((p_{1}-p_{2})\Delta V\)。按动能定理，这份功变成它的动能加势能增量：
\[
  (p_{1}-p_{2})\,\Delta V
  = \tfrac{1}{2}\rho\Delta V\,(v_{2}^{2}-v_{1}^{2})
  + \rho\Delta V\,g\,(h_{2}-h_{1}),
\]
两边约去 \(\Delta V\) 再移项，就是伯努利关系。看清楚了：它根本不是什么“流体的神秘新定律”，而是第 6 章的牛顿第二定律和第 9 章的能量守恒在流水上的联合记账。

力矩的完整写法带方向：\(\tau = rF\sin\theta\)，其中 \(\theta\) 是力臂方向与力的方向之间的夹角。\(\theta=0\) 时（力正对着支点拉）力矩为零，正是“推门轴门不转”的严格版；\(\theta=90^{\circ}\) 时力矩最大，所以扳手要垂直着拧。

固体的“弹簧账”也可以写成与形状无关的形式。定义应力 \(\sigma=F/S\)（单位面积分摊的内力）和应变 \(\varepsilon=\Delta L/L\)（相对伸长量），弹性范围内
\[
  \sigma = E\,\varepsilon,
\]
\(E\) 叫弹性模量，只由材料决定。钢的 \(E\) 约为 \(2\times10^{11}\) 帕：要让一根钢棒伸长百分之一，需要施加 \(2\times10^{9}\) 帕的应力——相当于在每平方厘米的截面上压两辆小汽车。这就是“固体都是弹簧，只是硬得离谱”的量化版本。橡皮筋的 \(E\) 只有钢的几万分之一，所以徒手可拉。
\end{mathbridge}

\begin{challenge}
伯努利关系有一条看不见的假设贯穿始终：忽略黏性。真实流体有“内摩擦”——蜂蜜从瓶口流出的速度远小于伯努利的预言，因为大量机械能沿途被黏性磨成了热。描述黏性是否重要的尺子叫雷诺数，它比较“惯性”和“黏性”谁说了算：雷诺数大（大船、快水、空气），惯性主导，伯努利近似好用；雷诺数小（微生物在体液里游动、细尘在静空气中沉降），黏性淹没一切，伯努利完全失效。一个细菌游泳时一停下划动就几乎瞬间停住，它活在“惯性可以忽略”的世界里——水对它就像蜂蜜对你一样黏稠。另一个深水区是机翼。严谨的空气动力学用“环量”描述机翼如何改变整个流场，再结合伯努利关系计算压强分布；而“机翼把气流向下偏转、气流反推机翼”是同一过程在动量语言里的描述。两种语言互不矛盾，矛盾只出现在把其中任何一种简化成“唯一的真相”。
\end{challenge}

\step{模型的边界}

\begin{mainline}
本章的每条公式都有清晰的国境线，越界必错。

\textbf{浮力公式的边界：需要引力，需要静平衡。}\(F_{\text{浮}}=\rho g V_{\text{排}}\) 里的 \(g\) 不是装饰——浮力的本质是“深处压强更大”，而压强随深度增加这件事本身就是引力拽着流体往下压出来的。在空间站的失重环境里，水不再分层，把乒乓球按到水下它也不会浮上来——那里根本没有“上”。这也是一个反直觉的检验：浮力不是水的“天生托举欲望”，而是引力场中的压强梯度。还有一个家用边界：物体必须让流体钻到自己底下才有浮力。桥墩深深插进河床淤泥，底面接触不到水，就\emph{不受浮力}；一个被吸盘紧紧压在泳池底的杯子，同样没有水从下面顶它。

\textbf{伯努利关系的边界：定常、不可压缩、黏性可忽略、沿同一流线。}四个条件缺一，账本就不平。电风扇出风口前后的气流、螺旋桨滑流，都不满足条件，不能直接套用。最典型的误用是飞机升力的民间解释：“机翼上表面路程长，空气要走完更长的路，所以流得快，伯努利说压强低，于是产生升力。”这段话每一环都经不起推敲——没有任何物理定律规定上下两股气流必须“同时到达”机翼后缘，实测中上方气流到达得\emph{更早}；按这个逻辑，上下对称的特技飞机机翼和倒飞动作根本无法产生升力，但它们都能飞。伯努利关系本身没有错，它忠实地把“流速分布”和“压强分布”互相翻译；但它不回答流速为什么那样分布。完整的图景是：机翼把经过的气流整体向下偏转，被向下推的空气反过来向上推机翼（牛顿第三定律），机翼上下表面的压强差与气流向下偏转是同一个过程的两面。请记住本书的原则：伯努利关系是升力账本上的重要一页，不是全部账本。

\textbf{帕斯卡原理的边界：流体近乎不可压缩。}在液压系统里压强传递近乎即时且不衰减；换成气体（可以大幅压缩），挤压的能量先被“存”进气体的压缩里，传递又慢又有损耗——所以刹车油管里一旦混进空气，刹车会变得软绵绵，这是真实的行车安全隐患。

\textbf{弹性模型的边界：弹性限度与疲劳。}胡克定律只在微小的相对形变内成立。越过弹性限度，材料发生永久变形（塑性）；继续加载则断裂。更阴险的是疲劳：反复弯折一根铁丝，每次的形变看似都在可恢复范围内，可几十次后它照样断——微观损伤在一次次的循环中积累。飞机机身每起降一次就经受一次加压减压循环，定期探伤检查就是和疲劳赛跑。
\end{mainline}

再集中列四条本章特有的典型误用：

\begin{itemize}[leftmargin=1.8em]
  \item \emph{“重的沉、轻的浮。”}错。裁判是物体所受重力与同体积流体所受重力之比（等价地，平均密度之比），不是物体自己的重量。
  \item \emph{“吸管靠嘴的吸力把饮料拉上来。”}错。嘴只是降低了管内的气压，把饮料压上来的是液面上的大气压。真空里“吸力”再强也喝不到。
  \item \emph{“流速快的地方压强低”随处乱套。}这条结论只在伯努利关系的四个条件下、沿同一流线成立。拿它解释一切“被风掀起”的现象，往往越界。
  \item \emph{“东西不动就是不受力。”}错。平衡是合力与合力矩同时为零；桥、塔、大坝内部时刻传递着巨大的内力，只是每一笔都被另一笔抵消。
\end{itemize}

\step{回头解释开场现象}

\textbf{倒扣水杯。}杯口面积约 \(40\) 平方厘米，大气压约 \(10^{5}\) 帕，空气从下面向上顶纸牌的力约为 \(10^{5}\times4\times10^{-3}\approx 400\) 牛——相当于四十千克重物的重力！杯中水不过半斤（约 \(3\) 牛），加上牌重，空气的大力士连零头都没用上。第一题的答案是（c），而且空气出手的量级是“几十千克重”。这个实验失败的常见方式恰好也是证据：杯口有缺口漏气、水没装满留下大气泡（气泡可被压缩，杯内压强建立不起足够的差值），牌就托不住了——失败方式恰好说明托牌的是压强差，不是“粘”或“吸”。

\textbf{台风掀屋顶。}近似用伯努利关系估算：风速取 \(50\) 米每秒，空气密度约每立方米 \(1.2\) 千克，\(\tfrac{1}{2}\rho v^{2}=\tfrac{1}{2}\times1.2\times2500=1500\) 帕。屋顶外侧气流快、压强低，屋内空气近乎静止、仍是正常大气压，于是一百平方米的屋顶受到的向上净力约 \(1500\times100=1.5\times10^{5}\) 牛——相当于十五吨重物的重力向上拔。这是真实数量级：台风天若能让门窗适度通风泄压，屋顶反而更安全（当然，真实屋面气流有分离和湍流，这只是估算）。吹纸条实验、两张纸相吸，都是同一个“快处压强低”的小号版本——你的纸条不是被吹起来的，是被下方静止空气顶进气流里的。

\textbf{双吸管与沉浮子。}两根吸管时，你嘴里的低压刚一建立，第二根吸管就把外面的大气放进嘴里补平，压强差无从谈起，液面上的大气压没有对手，饮料纹丝不动。沉浮子同理：捏瓶身时，帕斯卡原理把你加的压强传遍整瓶水，沉浮子内部的气泡被压缩、体积变小、排开的水变少，浮力敌不过重力，下沉；松手气泡复原，浮力恢复，上浮。潜水艇把水灌进压载舱变“重”下潜、排水上浮，用的是同一原理的另一端——它改的是自身重量，沉浮子改的是排水体积。

\textbf{桥、船、飞机的形状。}为什么三样东西长得完全不同？因为它们要用三种不同的方式回答同一个问题——把载荷安全地传到别处。桥：石头和混凝土抗压极强、抗拉极弱，所以拱桥把竖直的车辆载荷巧妙地转成沿拱身的\emph{压缩}，一直传进两岸山体；钢索抗拉极强，于是悬索桥反过来把载荷吊在受拉的主缆上；桁架桥则用一堆三角形杆件，把“弯”这种材料最怕的受力方式，拆解成每根杆单纯的拉或压。船：钢铁本身的密度是水的近八倍，唯一的活路是把身体做成空心，用尽量小的自重换来尽量大的排水体积——船壳的曲线就是浮力公式的几何化。飞机：机翼的弯度、迎角和面积，全部服务于一件事——在有限的速度下偏转足够多的空气、制造足够的上下压强差，同时还要控制阻力和结构重量。三种形状，三套权衡，同一条主线：结构和流体都在忠实地传递力，工程师的全部工作就是为力修一条又好又省的路。

\step{脑洞实验与挑战题}

\begin{thought}[把泡沫塑料杯沉入海底一万米]
太平洋马里亚纳海沟深约 \(11000\) 米。用 \(p=p_{0}+\rho g h\) 算沟底压强：\(\rho g h\approx1000\times10\times11000=1.1\times10^{8}\) 帕，约一千一百个大气压。深潜器的观察窗外，每平方厘米承受着约一吨的重量。现在开一个脑洞：把一只泡沫塑料咖啡杯用网兜挂在深潜器外面带下去（这是深海考察的真实传统），捞上来时它会缩成拇指大小、硬邦邦的纪念品——杯壁里无数小气泡中的气体被上千个大气压压没了。但请注意一个反直觉的细节：杯子的形状没有\emph{畸变}，只是均匀缩小，因为压强从四面八方同样地挤。均匀的高压本身并不压坏东西，压坏一样东西靠的是压强\emph{差}——深潜器的壳要厚，是因为它肚子里维持着一个大气压，壳外是一千一百个。同一次下潜里带下去的一块实心钢锭，捞上来几乎分毫未变：钢的体积模量极大，一千个大气压只能让它的体积缩小约百分之零点六。高压的世界比我们直觉里的“压扁”要讲道理得多。
\end{thought}

\begin{thought}[太空站里，油和水还分家吗]
在地面厨房里，油浮在水面上，界限分明——浮力公式说油密度小，被更重的水“挤”到上面。把这杯油水混合物带进失重的空间站，使劲摇匀后静置：一小时后再看，油滴仍然星星点点地悬在水中，迟迟不分层。原因很简单：浮力的前提 \(p=p_{0}+\rho g h\) 里的 \(g\) 失效了，流体内部不再有“越深压强越大”的梯度，重的成分没有任何理由沉底，轻的也没有理由上浮。同一场失重还会制造更多怪事：蜡烛火焰在地面上被热空气上升、冷空气补位的对流拉成泪滴形，在空间站里缩成一个安静的蓝色小球，而且因为新鲜空气送不进来，很快闷熄。这个脑洞一路通向后面的章节：压强、浮力这些“流体的脾气”究竟从哪来？第 27 章会把气体压强拆到分子层面——无数分子撞击器壁的统计平均；第 31 章则会告诉你，浮力、压强、温度这些宏观量，全都是微观大军的集体行为。本章的公式，是它们向宏观世界递交的简历。
\end{thought}

\begin{puzzles}
\item 一艘潜水艇悬浮在水中，既不上浮也不下沉，此时它排开水的重力与它自身的重力是什么关系？随后水兵向压载舱灌进一些海水，艇身外形几乎不变，接下来它会怎样？\emph{思路提示}：外形不变意味着 \(V_{\text{排}}\) 不变，浮力不变；变的是天平另一边的重力。判断沉浮，就是看这两个量的天平倒向哪边。
\item 两艘船在狭窄水道里并行，水手都知道要互相远离，否则两船会莫名其妙地“相吸”撞在一起。为什么？\emph{思路提示}：把两船之间看成一段被挤窄的水道，先用连续性方程判断那里的流速，再用伯努利关系比较两船“内侧”和“外侧”的压强。注意这个解释成立的条件：定常、近似无黏、沿流线——它在定性上是对的。
\item 古代石拱桥不用一根钉子、一根横梁，能把桥面举几百年；而同样材料搭成的平直石板桥，跨度稍大就会从中间断裂。石头怕什么、不怕什么？\emph{思路提示}：石头抗压极强、抗拉极弱。平直石板受载时，上沿被压、下沿被拉——断裂总从下沿开始。拱的曲线让载荷几乎全部沿拱身以\emph{压缩}的形式传向两岸。想想为什么悬索桥的材料选择恰好反过来。
\item 盒装牛奶只剪一个小口时，倒起来咕咚咕咚、时断时续；在盒子顶上再戳一个孔，奶就倒得又稳又顺。第二个孔没有改变奶的量，它改变了什么？\emph{思路提示}：奶流走后盒内气压会降低，盒外大气压就会顶着奶不让它出来。第二个孔是留给大气的“入场通道”。这和吸管、抽水机是同一出戏的第三幕。
\end{puzzles}

至此，我们从一张不落的纸牌出发，把“不动”拆成两笔账（力与力矩），把流体的推托拆成压强差，把沉浮交给“排开多少”，把流动交给能量守恒，把固体的硬与断交给藏在原子间的“弹簧”。下一章要换一个激进的视角：与其逐时刻追踪物体怎样运动，能不能把整条运动路径当成一个整体来比较，让“自然”自己在所有可能的路径中挑出真实的那一条？那将是从牛顿走向最小作用量原理的第一步。
